\chapter{Tests}

%--------------------------------------------------------------------------------

\unnumberedSection{Glossary Test}

This part introduces the key concepts of the Layout Control System and explains the major components that form the foundation of the system. The following summarizes the terminology and concepts used throughout this part of the book. It is intended as a quick reference and may safely be skipped on a first reading.

The LCS architecture is based on a distributed system where all nodes communicate via the layout control bus. All nodes receive broadcast messages and decide independently whether a message is relevant to them. Request/reply mechanisms are used where confirmation of an operation is required.

\begin{tsDefinition}

\tsDef{Node}
      {An addressable hardware module connected to the LCS bus.}

\tsDef{Node UID}
      {A globally unique hardware identifier assigned to a node.}

\tsDef{Port}
      {A functional unit within a node providing a specific capability.}

\tsDef{Channel}
      {An optional subdivision of a port providing multiple independent functions.}

\tsDef{Base Station}
      {The central coordinator responsible for managing locomotive sessions and system functions.}

\tsDef{Cab Handheld}
      {A user interface used by an operator to control locomotives.}

\tsDef{Locomotive Session}
      {A temporary association between a locomotive and the base station while the locomotive is under control.}

\tsDef{Event}
      {A broadcast notification used for communication between nodes.}

\tsDef{Attribute}
      {Configuration or status information belonging to a node or port.}

\tsDef{Node Identifier}
      {The logical address used to identify a node within the layout.}

\tsDef{Port Identifier}
      {The identifier of a port within a node.}

\tsDef{Channel Identifier}
      {The identifier of a channel within a port.}

\tsDef{\texttt{npId}}
      {The combined node, port, and channel identifier.}

\tsDef{\texttt{cabId}}
      {The logical identifier of a locomotive.}

\tsDef{\texttt{sId}}
      {The identifier of a locomotive session.}

\tsDef{\texttt{eventId}}
      {The unique identifier assigned to an event.}

\end{tsDefinition}

The following chapters introduce these concepts in more detail and explain how they are implemented in hardware, firmware, and software.

%--------------------------------------------------------------------------------
